\subsection{二次の近似}

微分係数による近似は
\[
 f(x)\simeq f(a)+f'(a)(x-a)
\]
という形であった.右辺は,点$a$で関数の値と傾きを合わせた一次式である.さらに二階微分の値も合わせたいなら,\,$(x-a)^2$の項を加える.\,$c(x-a)^2$を二回微分すると$2c$になるので,その係数は$f''(a)$そのものではなく,\,$f''(a)$の半分になる.同じ考えを続けると,\,$n$回までの微分の値を合わせた多項式は
\[
 P_n(x)=\sum_{k=0}^{n}\frac{f^{(k)}(a)}{k!}(x-a)^k
\]
となる.ここで$f^{(0)}=f$,\,$0!=1$とする.

この多項式が関数にどの程度近いかは,微分の値を合わせただけでは分からない.必要なのは,差$f(x)-P_n(x)$を調べる式である.まず微積分の基本定理から
\[
 f(x)=f(a)+\int_a^x f'(t)\dd t
\]
と書ける.積分の中で$f'(t)$を$f'(a)$に置き換えれば接線の式が現れるが,その置き換えで失われる分も残しておきたい.そのために,次の部分積分を行う.
\[
\begin{aligned}
 \int_a^x (x-t)f''(t)\dd t
 &=\bigl[(x-t)f'(t)\bigr]_a^x
   +\int_a^x f'(t)\dd t\\
 &=-(x-a)f'(a)+f(x)-f(a).
\end{aligned}
\]
したがって
\[
 f(x)=f(a)+f'(a)(x-a)
      +\int_a^x(x-t)f''(t)\dd t.
\]
最後の積分は,接線近似を正確な等式にするための項である.二階微分が小さい区間ではこの積分も小さくなる,という関係がここに現れている.

\begin{theorem}{有限テイラー公式と積分剰余}
$n$を$0$以上の整数とし,\,$f$は$a$と$x$を含む区間で$n+1$回連続微分可能であるとする.このとき
\[
 f(x)=\sum_{k=0}^{n}\frac{f^{(k)}(a)}{k!}(x-a)^k+R_n(x),
 \qquad
 R_n(x)=\frac{1}{n!}\int_a^x(x-t)^n f^{(n+1)}(t)\dd t
\]
が成り立つ.\,$R_n(x)$を剰余項という.
\end{theorem}

公式の仕組みを確かめる.\,$n=0$は微積分の基本定理そのものである.\,$n\geq1$のとき,剰余項を部分積分すると
\[
\begin{aligned}
 R_n(x)
 &=\left[\frac{(x-t)^n}{n!}f^{(n)}(t)\right]_a^x
   +\frac{1}{(n-1)!}\int_a^x(x-t)^{n-1}f^{(n)}(t)\dd t\\
 &=-\frac{f^{(n)}(a)}{n!}(x-a)^n+R_{n-1}(x).
\end{aligned}
\]
つまり,一つ前の剰余項から$n$次の項を取り出した残りが$R_n$である.この関係を順に使えば,公式が得られる.\,$x<a$でも,向きのある定積分の部分積分として同じ計算が成り立つ.

この導出で部分積分を選んだ理由も見ておこう.\,$(x-t)^n$を$t$で微分すれば次数が一つ下がり,\,$f^{(n+1)}(t)$を積分すれば微分の回数が一つ減る.その結果,一つ前の剰余項と結び付く.形式だけを覚えるよりも,この二つが同時に一段ずつ戻ることを覚えておくとよい.


